\section{Learning the complete binary qubit effect body}
\label{sec:commonlearn76}

We now pass from supplied descriptions to an unknown device.  Throughout
this section the input dimension is two.  Write
\[
 E=\tfrac12\bigl((1+b)I+r u\cdot\sigma\bigr),\qquad
 0\le r\le1,\qquad |b|+r\le1,
\]
with the direction absent when $r=0$.  The bias, contrast, and direction
are all unknown.  The outputs are ordered and classical, and each call
consumes its input without a residual quantum output.  A learning
procedure may prepare entangled input blocks, retain the unused input
qubits, and use classical feedback between blocks.  It has no access
to a measurement environment.  Its choices depend only on the public
parameters and the observed record.

The integer $N$ below is the horizon of a future comparison, measured
by $d_N$.  The number $M$ of training calls is a separate resource.
The output is a classical description of a legal effect.
For a unit vector $w$ set $P_w=(I+w\cdot\sigma)/2$ and write
$\mathcal M_{p,q,w}$ for the measurement associated with
$pP_w+q(I-P_w)$, where $p=(1+b+r)/2$ and $q=(1+b-r)/2$.
We use the spectral quantities
\[
 V=p(1-p)+q(1-q),\qquad
 K_N(p,q)=r\sqrt{N/(V+N^{-1})}.
\]

\begin{theorem}[Common learning at the finite-use scale]
\label{thm:commonlearn76}
There are absolute constants $c,C,\delta_0>0$ such that, for every
$N\ge1$, $0<\delta\le\delta_0$, and $0<\eta\le1/8$, one common
learning procedure returns an effect $\widehat E\in\mathfrak E_2$ with
\begin{equation}\label{eq:commonrisk76}
 \sup_{E\in\mathfrak E_2}
  \mathbb P_E\{d_N(\mathcal M_E,\mathcal M_{\widehat E})>\delta\}
       \le\eta.
\end{equation}
On every record it uses at most
\begin{equation}\label{eq:commoncalls76}
 M\le C N\delta^{-2}\log(1/\eta)
\end{equation}
calls, and every entangled input block has length at most $N$.
Conversely, every common procedure satisfying
\eqref{eq:commonrisk76}, with at most $M$ calls on every record,
obeys
\begin{equation}\label{eq:commonlower76}
 M\ge cN\delta^{-2}\log(1/\eta).
\end{equation}
The converse permits arbitrary reference-assisted adaptive training
and bounded public stopping.  The constants are uniform over the
closed effect body, including scalar, deterministic, and projective
effects.
\end{theorem}

The proof constructs one decision tree for the entire parameter body.
Its experiments are not selected from a known pair of hypotheses.
The angular part uses multiscale phase estimation, a method with
antecedents in \cite{HB76,KLY76}.  We give the phase and confidence
bounds explicitly.  A data-dependent amplitude threshold selects the
block length without prior calibration of the spectrum.

\subsection{Bernoulli errors and estimated axes}

For $s,t\in[0,1]$ put
\begin{equation}\label{eq:learnhell76}
 \mathfrak h(s,t)^2=(\sqrt s-\sqrt t)^2
                     +(\sqrt{1-s}-\sqrt{1-t})^2.
\end{equation}
Write $B_N(s,t)$ for the unhalved distance between the two
$N$-fold Bernoulli product laws.
We use the following elementary estimates at both support endpoints.

\begin{lemma}[Endpoint estimation]
\label{lem:learnendpoint76}
For every $s,t\in[0,1]$,
\begin{equation}\label{eq:learnproduct76}
 B_N(s,t)\le2\sqrt N\,\mathfrak h(s,t).
\end{equation}
For each fixed $A>0$ there is $C_A$ such that, if $0<\lambda\le1$ and
\begin{equation}\label{eq:learnvariance76}
 |s-t|\le A\bigl(\sqrt{t(1-t)\lambda}+\lambda\bigr),
\end{equation}
then $\mathfrak h(s,t)\le C_A\sqrt\lambda$.
If two estimates of an ordered pair $(p,q)$ each have
$\mathfrak h$ error at most $\varepsilon$, sorting the estimates
preserves this bound up to an absolute constant.
\end{lemma}
\begin{proof}
The root fidelity is $1-\mathfrak h^2/2$.  Tensorization and the
trace--fidelity inequality give
\[
 B_N(s,t)\le
 2\sqrt{1-(1-\mathfrak h(s,t)^2/2)^{2N}}
 \le2\sqrt N\,\mathfrak h(s,t).
\]
For the first square root in \eqref{eq:learnhell76}, if $t\ge\lambda$
divide the difference by $\sqrt t$ in
\eqref{eq:learnvariance76}.  If $t<\lambda$, both $t$ and $s$ are
at most $C_A\lambda$.  Apply the same argument to their complements.
For sorting, set $\phi(t)=\arcsin\sqrt t$.  Then
\[
 \mathfrak h(s,t)=2\sin\bigl(|\phi(s)-\phi(t)|/2\bigr).
\]
The two distances are uniformly comparable on $[0,\pi/2]$.
Sorting is nonexpansive for the sup norm against an already ordered
pair of real numbers.  Apply this fact in the $\phi$ coordinate.
\end{proof}

If $\widehat t$ is the empirical frequency in $n$ Bernoulli trials,
Bernstein's inequality gives, with failure at most $2e^{-L}$,
\begin{equation}\label{eq:learnbernstein76}
 |\widehat t-t|
 \le\sqrt{2t(1-t)L/n}+\frac{2L}{3n}.
\end{equation}
Thus Lemma~\ref{lem:learnendpoint76} applies with
$\lambda=L/n$, including at $t=0,1$.

For a unit estimate $w$ of $u$, write
\begin{equation}\label{eq:learnmisalignment76}
 h=\angle(u,w),\qquad
 \gamma=\tfrac r2(1-u\cdot w).
\end{equation}
The input states on the axes $+w$ and $-w$ have respective success
probabilities
\begin{equation}\label{eq:learnaxisprob76}
 s_+=p-\gamma,\qquad s_-=q+\gamma,
 \qquad p=\tfrac12(1+b+r),\quad q=\tfrac12(1+b-r).
\end{equation}
If $\widehat x$ estimates $x=ru$ with Euclidean error at most $e$,
normalizing $\widehat x$ gives
\begin{equation}\label{eq:learnaxiserror76}
 rh\le\pi\min(r,e),\qquad
 \gamma\le\min(r,e^2/r)\quad(r>0).
\end{equation}
Indeed $\|u-w\|\le2e/r$, $h\le(\pi/2)\|u-w\|$, and
$\gamma=r\|u-w\|^2/4$.  If $\widehat x=0$, choose a fixed direction;
then $e\ge r$ gives the same bounds.  At $r=0$ both quantities vanish.

\begin{lemma}[Product learning away from unit contrast]
\label{lem:productlearn76}
Suppose $r\le7/8$.  There is a product-input procedure, not requiring
$b,r,u$ as input, which uses eight groups of $n$ calls and returns a
legal effect $F$ such that, for $1\le L\le n$,
\begin{equation}\label{eq:productlearn76}
 d_N(\mathcal M_E,\mathcal M_F)
       \le C\bigl(\sqrt{NL/n}+NL/n\bigr)
\end{equation}
outside an event of probability at most $Ce^{-L}$.
\end{lemma}
\begin{proof}
First use the positive and negative eigenstates of each of the three
Pauli axes, with $n$ calls per state.  Encode an observed outcome as
$Z\in\{-1,1\}$.  The two means on axis $j$ are
$b+x_j$ and $b-x_j$.  Estimate $x_j$ by half their empirical
difference.  Put
\[
 w_0=1-b^2,\qquad \lambda=L/n.
\]
The two variances on axis $j$ sum to $2(w_0-x_j^2)$.
Bernstein's inequality and a union bound therefore give
\begin{equation}\label{eq:learnproductvector76}
 e=\|\widehat x-x\|
       \le3\bigl(\sqrt{w_0\lambda}+\lambda\bigr).
\end{equation}
More explicitly, the coordinate error is at most
$\sqrt{2w_0\lambda}+4\lambda/3$; the sign observation has range
length two.  Taking the Euclidean norm gives
\eqref{eq:learnproductvector76}, since $\sqrt6<3$ and
$4\sqrt3/3<3$.

Normalize to a direction $w$.  Feasibility and $r\le7/8$ imply
\begin{equation}\label{eq:learnlowvariance76}
 w_0\ge r(2-r),\qquad
 V=\tfrac12(w_0-r^2)\ge w_0/9.
\end{equation}
Use \eqref{eq:learnaxiserror76} and
Theorem~\ref{thm:biasedangular75}.  The square-root term in
\eqref{eq:learnproductvector76} contributes at most
$C\sqrt{N\lambda}$ to $K_N(p,q)h$, and its additive term
contributes at most $CN\lambda$.  Consequently,
\begin{equation}\label{eq:learnlowangle76}
 d_N(\mathcal M_{p,q,u},\mathcal M_{p,q,w})
       \le C\bigl(\sqrt{N\lambda}+N\lambda\bigr).
\end{equation}
This calculation does not require a lower bound on $r$ or $w_0$.

It remains to control both endpoint probabilities in
\eqref{eq:learnaxisprob76}.  For either $t=p$ or $t=q$ put
$v_t=t(1-t)$.  Since
\[
 |p(1-p)-q(1-q)|=r|b|\le r,
 \qquad w_0=2V+r^2\le4v_t+3r,
\]
we claim
\begin{equation}\label{eq:learngammabound76}
 \gamma\le9\sqrt{v_t\lambda}+72\lambda.
\end{equation}
Equation~\eqref{eq:learnproductvector76} gives
$e^2\le18(w_0\lambda+\lambda^2)$.  If $r\le\lambda$, use
$\gamma\le r$.  If $r>\lambda$, use
\[
 \gamma\le72v_t\lambda/r+72\lambda,
 \qquad \gamma\le r,
\]
and $\min(r,72v_t\lambda/r)\le\sqrt{72v_t\lambda}$.
This proves \eqref{eq:learngammabound76}.

Make two fresh groups of $n$ calls, along $+w$ and $-w$.
Conditional on the previous record, these are Bernoulli trials with
the parameters in \eqref{eq:learnaxisprob76}.  Their variances differ
from $v_t$ by at most $\gamma$.  Equations
\eqref{eq:learnbernstein76} and \eqref{eq:learngammabound76} imply
\[
 |\widehat t_{\rm raw}-t|
       \le C\bigl(\sqrt{v_t\lambda}+\lambda\bigr).
\]
For the additional term use
$\sqrt{\gamma\lambda}
 \le C(\sqrt{v_t\lambda}+\lambda)$, which follows from
$\sqrt{AB}\le(A+B)/2$ with
$A=\sqrt{v_t\lambda}$ and $B=\lambda$.
Lemma~\ref{lem:learnendpoint76} gives root-distance error
$C\sqrt\lambda$ for both endpoints.  Sort the two empirical
parameters to obtain $\widehat p\ge\widehat q$ and return
$F=\widehat pP_w+\widehat q(I-P_w)$.  The aligned programme
bound \eqref{eq:alignedspectral75} and
\eqref{eq:learnproduct76} give a spectral contribution at most
$C\sqrt{N\lambda}$.  Combine with \eqref{eq:learnlowangle76}.
All concentration bounds for the fresh groups hold conditionally,
so their failure probabilities may be added to those of the first
six groups.
\end{proof}

\subsection{A common entangled experiment with unknown bias}

\begin{lemma}[Two-plane phase observation]
\label{lem:ghzreadout76}
Fix an orthonormal frame $(e_1,e_2,w)$ known to the experimenter.
For each $i=1,2$ and integer $m\ge1$, there is an $m$-call experiment
with an observation $Y\in\{-1,1\}$ whose mean, for a chosen phase
$\varphi$, is
\begin{equation}\label{eq:ghzreadout76}
 \mathbb E_EY=\Re(e^{i\varphi}Z_i),\qquad
 Z_i=\bigl[x\cdot w+i x\cdot e_i\bigr]^m,
 \quad x=ru.
\end{equation}
Two phases estimate the complex number $Z_i$.  For $0<\zeta\le1$
and $0<\epsilon\le1/2$, four settings, using the two planes,
estimate both complex numbers with error at most $\zeta$ and
failure at most $\epsilon$, with at most
$C\zeta^{-2}\log(C/\epsilon)$ independent blocks per setting.

Fix $a=1/64$.  If $r>0$ and $h=\angle(u,w)\le2a/m$, then
\begin{equation}\label{eq:ghzamplitude76}
 .99r^m\le |Z_i|\le r^m,
 \qquad
 \operatorname{Arg}Z_i
       =m\arctan\frac{u\cdot e_i}{u\cdot w},
\end{equation}
where the principal argument is used.  If the amplitudes are bounded
below by a fixed positive constant and $\zeta$ is sufficiently small,
these observations produce a direction estimate $w'$ with
\begin{equation}\label{eq:ghzphaseerror76}
 \angle(u,w')\le C\zeta/m.
\end{equation}
\end{lemma}
\begin{proof}
Choose a computational basis along the axis perpendicular to the
plane spanned by $w,e_i$, with phase convention such that the
off-diagonal matrix element of
$O=2E-I=bI+ru\cdot\sigma$ is
$r(u\cdot w+i u\cdot e_i)$.  Prepare
\[
 |G_{\varphi,s}\rangle
       =\frac{|0\rangle^{\otimes m}
                    +s e^{i\varphi}|1\rangle^{\otimes m}}{\sqrt2},
\]
where $s$ is a fresh fair sign.  Pass its qubits through the
measurement and record $Y=s\prod_{j=1}^m Z_j$, where each output
is encoded as $Z_j\in\{-1,1\}$.  Memorylessness makes the joint
outcome law the product POVM on the entangled input; the qubits may
be supplied sequentially while the unused ones are retained.

The expansion of $\langle G_{\varphi,s}|O^{\otimes m}
|G_{\varphi,s}\rangle$ has diagonal term
\[
 \tfrac12\bigl((b+ru\cdot f_i)^m+(b-ru\cdot f_i)^m\bigr),
\]
where $f_i$ is the chosen perpendicular axis.  It is independent
of $s$ and disappears after multiplication by $s$ and averaging.
The off-diagonal terms give \eqref{eq:ghzreadout76}.  At $x=0$
the mean is zero and no direction is needed.  In particular
the experiment removes the unknown bias exactly.  The phases
$0,\pi/2$ give the real part and the negative imaginary part.
Hoeffding's inequality for four bounded sample means proves the
stated block count.

For the phase conclusions assume $r>0$ and put
$\theta_i=\arctan((u\cdot e_i)/(u\cdot w))$.  Then
$|m\theta_i|\le2a$ and
\[
 |Z_i|=r^m\bigl(1-(u\cdot f_i)^2\bigr)^{m/2}
       \ge r^m\cos(h)^m
       \ge r^m(1-2a^2/m).
\]
This proves \eqref{eq:ghzamplitude76}, with no phase ambiguity.
For the estimates set
\[
 \widehat\theta_i=\operatorname{Arg}(\widehat Z_i)/m,
 \qquad
 w'=\frac{w+\tan(\widehat\theta_1)e_1
                   +\tan(\widehat\theta_2)e_2}
              {\|w+\tan(\widehat\theta_1)e_1
                   +\tan(\widehat\theta_2)e_2\|}.
\]
The exact formula with $\theta_i$ recovers $u$.  For
$|\widehat Z_i-Z_i|\le|Z_i|/2$,
\[
 |\operatorname{Arg}(\widehat Z_i)-\operatorname{Arg}(Z_i)|
       \le2|\widehat Z_i-Z_i|/|Z_i|.
\]
On $|\theta_i|\le1/4$ the displayed map from two angles to a unit
vector is Lipschitz; a bound of $16$ for angular distance against
the largest coordinate error follows from differentiation of
$\tan$ and normalization.  This proves
\eqref{eq:ghzphaseerror76}.  Our fixed cone lies strictly inside
the principal-argument interval, and the argument uses no
unwrapping threshold from a separate phase-estimation theorem.
\end{proof}

\begin{lemma}[Selecting a block length from the observations]
\label{lem:noiseblock76}
Suppose $r\ge5/8$ and an initial known direction $w$ satisfies
$\angle(u,w)\le a=1/64$.  A common adaptive procedure selects a
length $m_*$ and a direction $w_*$, for $N\ge1$ and
$0<\eta\le1/8$, using at most $CN\log(C/\eta)$ calls on every
record.  With failure at most
$\eta/8$,
\begin{align}
 \tfrac18\min\{N,(1-r)^{-1}\}&\le m_*\le N,
                                      \label{eq:learnblockscale76}\\
 r^{m_*}&\ge1/5,
 &\angle(u,w_*)&\le a/m_*.
                                      \label{eq:learnblockaxis76}
\end{align}
The reciprocal at $r=1$ is interpreted as infinity.  New observations
at $m_*,w_*$ have both complex amplitudes at least $1/6$.
\end{lemma}
\begin{proof}
Fix $\kappa=a/1024$, and let $H$ be the largest power of two not
exceeding $N$.  Try lengths $m=1,2,4,\ldots,H$.  At length $m$,
estimate both complex numbers in Lemma~\ref{lem:ghzreadout76} to
error at most $\kappa$, assigning conditional failure allowance
\begin{equation}\label{eq:learnstagebudget76}
 \epsilon_m=\frac{\eta m}{32N}.
\end{equation}
Accept the length if both empirical amplitudes are at least $1/4$.
On acceptance update the direction by the formula in the proof of
Lemma~\ref{lem:ghzreadout76}, and double the length unless the cap
has been reached.  On rejection retain the preceding accepted
length and direction.  If none was accepted, use a prescribed
fallback with $m_*=1$.  All choices are functions of the record.

We prove the assertions on the event that the successive complex
estimates obey their bounds.  Before a length $m$ is tried the
angular error is at most $2a/m$: this holds initially, and is
preserved by each accepted update.  On acceptance the true
amplitudes are at least $1/4-\kappa>1/5$.  The argument bound and
the Lipschitz constant in Lemma~\ref{lem:ghzreadout76} give an
updated angular error at most
\[
 \frac{32\kappa}{(1/4-\kappa)m}\le\frac{a}{4m}.
\]
Furthermore $r^m\ge1/4-\kappa$, whence $m(1-r)\le2$.

On rejection, \eqref{eq:ghzamplitude76} gives $r^m<1/3$.
Since $r\ge5/8$,
\[
 -\log r\le(1-r)/r\le2(1-r),
 \qquad m(1-r)\ge(\log3)/2.
\]
The preceding accepted length $m/2$ is therefore at least
$(\log3)/(4(1-r))$.  Length one is always accepted on the good
event, since its true amplitudes exceed $.99(5/8)$.  If the cap is
reached, the selected length is $H\ge N/2$.  These facts prove
\eqref{eq:learnblockscale76} and \eqref{eq:learnblockaxis76}.
At the updated direction the amplitude is at least
$r^{m_*}\cos(a/m_*)^{m_*}>1/6$.

The concentration bounds are conditional on each preceding good
record.  Their failure probabilities sum to at most $\eta/8$,
because the sum of the tested powers of two is at most $2N$.
By Lemma~\ref{lem:ghzreadout76}, the number of calls is bounded by
an absolute constant times
\begin{equation}\label{eq:learndyadiccost76}
 \sum_{m=1,2,4,\ldots,H}
       m\log\frac{CN}{\eta m}
       \le C'N\log(C'/\eta).
\end{equation}
To see this, write $m=H/2^k$ and use the convergent sums of
$2^{-k}$ and $k2^{-k}$, together with $1\le N/H<2$.
The public cap and these budgets bound the cost on every record,
including records on which an estimate fails.  Rejection can only
shorten the sum.
\end{proof}

\subsection{The learning bound and its converse}

\begin{proof}[Proof of Theorem~\ref{thm:commonlearn76}]
We choose constants in the following order.  Fix the local cone
$a=1/64$ and $\kappa=a/1024$ as above.  Choose a sufficiently small
absolute constant $c_*>0$ in the final phase tolerance below.
Then choose the absolute multipliers in the Bernoulli sample counts
sufficiently large.  Finally reduce $\delta_0$, if needed, so that
$\delta_0\le1/4096$.  No optimized value of these constants is used.

First make the six Pauli-axis product experiments, with a sufficiently
large absolute multiple of $\log(C/\eta)$ calls per input.  With
failure at most $\eta/8$, the half-difference estimate satisfies
\[
 \|\widehat x-x\|\le\varepsilon_0,
 \qquad \varepsilon_0\le a/100.
\]
If $\|\widehat x\|<3/4$, choose the product procedure.  On the good
event $r\le7/8$, so Lemma~\ref{lem:productlearn76}, with
\[
 L=\log(C/\eta),\qquad
 n=\left\lceil CN\delta^{-2}L\right\rceil,
\]
gives error at most $\delta/2$ with the asserted call count, after
increasing $C$.

Otherwise $r\ge5/8$, and the normalized initial estimate has
angular error at most $a$.  Apply Lemma~\ref{lem:noiseblock76}.
At the selected length and updated direction, make fresh
two-plane observations with complex error tolerance
\begin{equation}\label{eq:learnfinetolerance76}
 \zeta_*=c_*\delta\sqrt{m_*/N},
\end{equation}
and failure allowance $\eta/8$.  The resulting direction $w_f$
satisfies
\begin{equation}\label{eq:learnfineaxis76}
 h_f=\angle(u,w_f)
       \le Cc_*\delta/\sqrt{Nm_*}.
\end{equation}
The calls used in this step are at most
\begin{equation}\label{eq:learnfinecost76}
 Cm_*\zeta_*^{-2}\log(C/\eta)
       \le C'N\delta^{-2}\log(C/\eta).
\end{equation}

Put $\nu=1-r$.  Since $|b|\le\nu$ and $r\ge5/8$,
\[
 r\nu\le V=\tfrac12(1-r^2-b^2)\le\nu.
\]
Equation~\eqref{eq:learnblockscale76} implies
\begin{equation}\label{eq:learnblockvariance76}
 m_*(V+N^{-1})\ge c_0>0
\end{equation}
for an absolute $c_0$.  Thus \eqref{eq:learnfineaxis76} and
Theorem~\ref{thm:biasedangular75} give
\begin{equation}\label{eq:learnhighangle76}
 d_N(\mathcal M_{p,q,u},\mathcal M_{p,q,w_f})
       \le Cc_*\delta.
\end{equation}
The endpoint misalignment is also bounded:
\begin{equation}\label{eq:learnhighgamma76}
 \gamma=\tfrac r2(1-u\cdot w_f)
       \le rh_f^2/4\le Cc_*^2\delta^2/N.
\end{equation}

Make fresh groups of $n=\lceil CN\delta^{-2}\log(C/\eta)\rceil$
calls along $+w_f$ and $-w_f$.  Conditional Bernstein bounds,
Lemma~\ref{lem:learnendpoint76}, and
$\mathfrak h(s,t)^2\le2|s-t|$ give
\begin{align}
 \mathfrak h(\widehat p_{\rm raw},p),\,
 \mathfrak h(\widehat q_{\rm raw},q)
 &\le C\sqrt{\log(C/\eta)/n}+C\sqrt\gamma\notag\\
 &\le c'\delta/\sqrt N.
                                      \label{eq:learnhighend76}
\end{align}
The constant $c'$ can be made as small as required by the preceding
choices of $c_*$ and the sample multiplier.  Sort the estimates and
return $\widehat E=\widehat pP_{w_f}+\widehat q(I-P_{w_f})$.
Equations~\eqref{eq:alignedspectral75} and
\eqref{eq:learnproduct76} make the spectral loss at most
$\delta/4$; choose $c_*$ so that
\eqref{eq:learnhighangle76} contributes at most $\delta/4$.
The total failure probability is less than $\eta$.

On exceptional records fix arbitrary legal conventions for a zero
complex estimate, an ambiguous principal argument, an estimated
angle outside $[-1/4,1/4]$, or an absent accepted length.  Each
output remains a legal effect and each block
length remains at most $N$.  The public budgets in
\eqref{eq:learndyadiccost76}, \eqref{eq:learnfinecost76}, and the
product experiments give \eqref{eq:commoncalls76} on every record.
This proves the upper bound, with error slack.

For the converse consider the scalar effects
\begin{equation}\label{eq:learnscalarpair76}
 E_0=\tfrac12I,
 \qquad E_1=(\tfrac12+\Delta)I,
 \qquad \Delta=256\delta/\sqrt N.
\end{equation}
Since $\delta\le1/4096$, both parameters lie in $[1/2,9/16]$.
The finite Bernoulli lower bound in
Lemma~\ref{lem:bernoulliproduct75} yields
\begin{align}
 d_N(\mathcal M_{E_0},\mathcal M_{E_1})
 &=B_N(1/2,1/2+\Delta)\notag\\
 &\ge\tfrac1{64}
       \min\{1,\Delta\sqrt{N/(5/4)}\}>2\delta.
                                      \label{eq:learnscalarsep76}
\end{align}
For these devices every call gives a fresh Bernoulli bit independent
of the input.  On an entangled input, each conditional reference
block is the same reduced state multiplied by the scalar outcome
probability.  Thus any adaptive learning record, including its
final classical estimate, is a common stochastic processing of
$M$ independent Bernoulli bits and parameter-independent resources.
Publicly stopped records may be padded.  The relative entropy before
processing is at most
\begin{equation}\label{eq:learnscalarkl76}
 M D\bigl(\operatorname{Bern}(1/2)\,\|\,
                \operatorname{Bern}(1/2+\Delta)\bigr)
   =-\tfrac M2\log(1-4\Delta^2)
   \le3M\Delta^2.
\end{equation}
The two radius-$\delta$ success balls are disjoint by
\eqref{eq:learnscalarsep76}.  Testing whether the estimated effect
belongs to the first ball therefore gives two errors at most $\eta$.
For classical laws $P,Q$, Jensen's inequality gives root fidelity
at least $e^{-D(P\|Q)/2}$.  The trace--fidelity inequality then
implies that the sum of testing errors is at least
$\tfrac12e^{-D(P\|Q)}$.  Consequently
\[
 2\eta\ge\tfrac12e^{-3M\Delta^2},
 \qquad
 M\ge\frac{N}{3\cdot256^2\delta^2}
                    \log\frac1{4\eta}.
\]
For $\eta\le1/8$ this proves \eqref{eq:commonlower76}.
\end{proof}

\begin{corollary}[A learned joint description]
\label{cor:learncode76}
For rational $\delta$ in the range of
Theorem~\ref{thm:commonlearn76}, the learner may return one word of
the rational joint code, with the same risk and call-count bounds.
Its payload length is
\begin{equation}\label{eq:learnpayload76}
 2\log_2N+\log_2\log(N+2)
             +4\log_2(1/\delta)+O(1),
 \qquad 0<\delta\le\delta_0.
\end{equation}
The decoding centre is a legal rational effect.  This is a statement
about device calls and payload length; classical arithmetic,
workspace, and state-preparation costs are separate resources.
\end{corollary}
\begin{proof}
Use the learner with loss budget $\delta/2$.  Its output is a finite
computable expression determined by a finite record.  Choose a
positive rational $\varepsilon\le\delta/(100N)$ and rational
approximations $(b_0,x_0)$ to its bias and Bloch vector with
\[
 |b_0-\widehat b|+\|x_0-\widehat x\|\le\varepsilon.
\]
Set $(b_1,x_1)=(b_0,x_0)/(1+2\varepsilon)$.  Then
\[
 |b_1|+\|x_1\|
       \le\frac{1+\varepsilon}{1+2\varepsilon}\le1,
\]
and the exact one-use formula \eqref{eq:biasedoneuse75} bounds the
distance from the unrounded hypothesis by $3\varepsilon$.
The $N$-use hybrid inequality adds at most $3N\varepsilon$.
Apply Theorem~\ref{thm:biasedcodec75} to this rational estimated
effect with public accuracy $\delta/2$.  Its certified radius is
$\delta/4$, so the total loss is less than $\delta$.
The capacity estimate in that theorem, at a constant rescaling of
accuracy, gives \eqref{eq:learnpayload76}.  This postprocessing uses
no additional device calls and receives no exact description of the
unknown target.
\end{proof}

The loss in Theorem~\ref{thm:commonlearn76} is a future adaptive
distance.  One-use measurement tomography is studied in
\cite{MB67,ZRK76}; combining the stated one-use accuracy bounds
only with the hybrid inequality gives a quadratic dependence on
the future horizon.  The argument here uses finite-use geometry
and data-dependent entangled blocks to obtain the linear dependence
in \eqref{eq:commoncalls76}, with a matching scalar converse.
Learning a stored programme from which a measurement is later
retrieved, as in \cite{BDPS76}, has a different output resource.
The present theorem concerns a classical estimate over the entire
ordered binary qubit body; it does not assert this learning bound
for the arbitrary-dimensional matrix body of
Theorem~\ref{thm:matrixmetric76}.
